\chapter{Modélisation et conception}
\thispagestyle{MyStyle}
\begin{onehalfspace}

\section{Introduction}

Ce chapitre présente la conception de la plateforme de gestion des candidatures. Après l'expression des besoins, il s'agit de traduire ces spécifications en modèles UML, afin de clarifier le comportement du système et sa structure avant de présenter, dans le chapitre suivant, les outils et les captures d'écran de la solution réalisée.

Nous commençons par justifier le choix d'UML, puis nous présentons les diagrammes de cas d'utilisation, qui identifient les acteurs et les fonctionnalités. Les diagrammes de séquence décrivent ensuite les interactions pour les principaux scénarios métier. Enfin, les diagrammes de classes formalisent les entités et leurs relations. Cette démarche vise à assurer la cohérence entre les besoins exprimés au chapitre~2 et la solution conçue.

\section{Modélisation}

\subsection{Le choix d'UML}

La conception du système a été formalisée à l'aide d'UML (Unified Modeling Language)~\cite{uml}, langage de modélisation normalisé, largement utilisé pour décrire le comportement et la structure d'une application avant sa réalisation. Ce choix s'inscrit dans la démarche en cascade retenue pour le projet : les besoins du chapitre~2 sont d'abord traduits en modèles, validés avec l'encadrant, puis servent de référence au développement.

UML permet de représenter le système selon plusieurs points de vue complémentaires, sans ambiguïté excessive et de manière indépendante des détails d'implémentation. Trois types de diagrammes ont été retenus, car ils couvrent l'essentiel de la conception métier :

\begin{itemize}
  \item les diagrammes de cas d'utilisation, qui identifient les acteurs et les fonctionnalités accessibles ;
  \item les diagrammes de séquence, qui décrivent le déroulement des principaux scénarios (messages échangés et résultat attendu) ;
  \item les diagrammes de classes, qui structurent les entités du domaine et leurs relations.
\end{itemize}

Cette modélisation vise à garantir la cohérence entre les besoins exprimés (RF-01 à RF-34) et la solution conçue, en particulier sur la séparation des rôles, l'enchaînement import--répartition--convocations, et les règles d'affectation des candidats.

\subsection{Outil de modélisation : PlantUML}

\begin{figure}[H]
  \centering
  
\includegraphics[width=0.32\textwidth]{plantuml-logo.jpg}
  \caption{Logo de PlantUML}
  \label{fig:logo-plantuml}
\end{figure}

PlantUML~\cite{plantuml} a été utilisé pour produire les diagrammes de modélisation présentés dans ce chapitre, notamment les diagrammes de cas d'utilisation, de séquence et de classes. Son approche textuelle facilite la mise à jour des diagrammes et leur intégration dans le cycle de documentation.

Ce choix permet de conserver des diagrammes versionnables avec le projet, contrairement à des outils purement graphiques dont les modifications sont plus difficiles à suivre dans un dépôt Git.

\section{Analyse et conception générale}

\subsection{Diagrammes de cas d'utilisation}

Les diagrammes de cas d'utilisation offrent une vue fonctionnelle du système. Ils précisent qui interagit avec l'application et quelles actions sont autorisées, sans décrire encore le détail des traitements. Le diagramme général ci-dessous rassemble l'ensemble des cas identifiés pour la plateforme de dématérialisation de la gestion des candidatures.

\paragraph{Identification des acteurs}

Deux acteurs humains apparaissent sur le diagramme. Ils correspondent aux profils définis dans le cahier des charges.

\begin{table}[H]
\centering
\caption{Acteurs du diagramme de cas d'utilisation général}
\label{tab:uc-acteurs}
\begin{spacing}{1}
\resizebox{\textwidth}{!}{%
\begin{tabular}{|p{3.4cm}|p{12.0cm}|}
\hline
\rowcolor[HTML]{C9DAF8}
\textbf{Acteur} & \textbf{Rôle vis-à-vis du système} \\
\hline
Gestionnaire &
Acteur métier principal. Il consulte le tableau de bord et les listes (candidats, concours, lieux, répartition, convocations). Il est le seul à gérer les candidats, les concours, les centres, établissements et lieux, à lancer la répartition et à gérer les convocations. \\
\hline
Administrateur &
Il consulte les mêmes informations métier en lecture. Il est le seul à gérer les comptes. Il ne lance pas la répartition et n'assure aucune opération de gestion métier. \\
\hline
\end{tabular}%
}
\end{spacing}
\end{table}

Le candidat n'apparaît pas sur ce diagramme : il ne dispose pas d'espace dans l'application. Il reste un destinataire externe des convocations.

\paragraph{Diagramme de cas d'utilisation général}

\begin{figure}[H]
  \centering
  \includegraphics[width=\textwidth,height=0.85\textheight,keepaspectratio]{usecase-general.png}
  \caption{Diagramme de cas d'utilisation général}
  \label{fig:usecase-general}
\end{figure}

Le diagramme de la figure~\ref{fig:usecase-general} situe le système dans son ensemble, avec deux acteurs~: le gestionnaire, à gauche, et l'administrateur, à droite. Tous les cas d'utilisation incluent le cas \textit{S'authentifier}~: aucune fonctionnalité n'est accessible sans connexion préalable.

Les cas de consultation --- tableau de bord, listes des candidats, des concours, des lieux d'examen, de la répartition et des convocations --- sont partagés par les deux acteurs. Les opérations d'écriture métier (gestion des candidats, des concours, des centres, établissements et salles, lancement de la répartition, gestion des convocations) sont en revanche réservées au gestionnaire. L'administrateur est le seul à \textit{Gérer les comptes}~; il n'intervient pas sur les données métier et ne déclenche pas la répartition.

Cette organisation traduit la règle posée au chapitre~2~: le gestionnaire pilote le processus opérationnel, l'administrateur supervise en lecture et administre les accès.

\paragraph{Description des cas d'utilisation}

\begin{table}[H]
\centering
\caption{Description des cas d'utilisation}
\label{tab:uc-description}
\begin{spacing}{1}
\resizebox{\textwidth}{!}{%
\begin{tabular}{|p{4.4cm}|p{4.0cm}|p{6.8cm}|}
\hline
\rowcolor[HTML]{C9DAF8}
\textbf{Cas d'utilisation} & \textbf{Acteur(s)} & \textbf{Description} \\
\hline
S'authentifier &
Gestionnaire, Administrateur &
Cas inclus par tous les autres. Il conditionne l'accès au système. \\
\hline
Consulter le tableau de bord &
Gestionnaire, Administrateur &
Accès aux indicateurs de suivi (effectifs, remplissage, activité récente). \\
\hline
Consulter la liste des candidats &
Gestionnaire, Administrateur &
Lecture de la liste des candidats importés. \\
\hline
Consulter la liste des concours &
Gestionnaire, Administrateur &
Lecture du référentiel des concours. \\
\hline
Consulter la liste des centres, établissements et salles &
Gestionnaire, Administrateur &
Lecture du référentiel des lieux d'examen. \\
\hline
Consulter la répartition &
Gestionnaire, Administrateur &
Lecture des affectations (centre, établissement, salle, place). \\
\hline
Consulter la liste des convocations &
Gestionnaire, Administrateur &
Lecture des convocations générées. \\
\hline
Gérer les candidats &
Gestionnaire &
Import Excel, modification, suppression et, le cas échéant, affectation manuelle. \\
\hline
Gérer les concours &
Gestionnaire &
Création, modification et suppression des concours. \\
\hline
Gérer centres, établissements et lieux &
Gestionnaire &
Création, modification et suppression des centres, établissements et salles. \\
\hline
Lancer la répartition &
Gestionnaire &
Déclenchement de la répartition automatique des candidats dans les salles. \\
\hline
Gérer les convocations &
Gestionnaire &
Génération des convocations et envoi groupé par courrier électronique. \\
\hline
Gérer les comptes &
Administrateur &
Création, modification et suppression des comptes gestionnaires. \\
\hline
\end{tabular}%
}
\end{spacing}
\end{table}

\subsection{Diagrammes de séquence}

Les diagrammes de séquence décrivent le déroulement temporel des principaux scénarios. Ils précisent les messages échangés entre l'application web, la passerelle API, les microservices et PostgreSQL.

\subsubsection{Diagramme de séquence --- authentification}

\fullpagefig{seq-authentification.png}{Diagramme de séquence --- authentification}{fig:seq-auth}

Le diagramme concerne l'utilisateur --- gestionnaire ou administrateur --- qui accède à la page de connexion et saisit son nom d'utilisateur et son mot de passe. L'application web vérifie que les champs sont renseignés, puis envoie une requête \texttt{POST /auth/login} à la passerelle API, qui la transmet au service d'authentification. Celui-ci interroge PostgreSQL, puis contrôle que le compte est actif et que le mot de passe est correct. En cas d'échec (identifiants invalides ou compte désactivé), un \texttt{401 Unauthorized} revient par la passerelle et un message d'erreur s'affiche. En cas de succès, le service génère un JWT d'accès (nom d'utilisateur et rôle), répond \texttt{200 OK}, l'application stocke le jeton et ouvre l'accès au tableau de bord. Les requêtes suivantes passent de la même façon par la passerelle, avec l'en-tête \texttt{Authorization: Bearer}.

\subsubsection{Diagramme de séquence --- gestion des candidats}

Le diagramme de la figure~\ref{fig:seq-candidats} décrit les opérations du gestionnaire sur la page «~Candidats~». Cinq participants interviennent~: le gestionnaire, l'application web (React), la passerelle API, le service candidats (Spring Boot) et PostgreSQL. Toutes les requêtes transitent par la passerelle, avec l'en-tête \texttt{Authorization: Bearer}. Il n'existe pas de création unitaire~: les candidats sont créés uniquement par import Excel.

\fullpagefig{seq-candidats.png}{Diagramme de séquence --- gestion des candidatures}{fig:seq-candidats}

Le scénario se déroule en cinq temps, correspondant aux cadres du diagramme.

Consulter la liste. À l'ouverture de la page, l'application web demande la liste par \texttt{GET /api/candidats}. La passerelle transmet la requête au service candidats, qui exécute un \texttt{SELECT} sur PostgreSQL. La réponse \texttt{200 OK} ramène la liste, ensuite affichée à l'écran.

Importer un fichier Excel. Le gestionnaire choisit un fichier \texttt{.xlsx} des candidats admis. L'application contrôle l'extension, puis envoie le fichier en \texttt{POST /api/candidats/import}. Si le fichier est vide ou d'un autre format, le service répond \texttt{400 Bad Request} et un message d'erreur s'affiche. Sinon, le service lit le tableur, valide les lignes, résout le concours, puis insère ou met à jour chaque candidat selon le CIN. La réponse \texttt{200 OK} indique les volumes traités (insérés, mis à jour, ignorés, erreurs). L'application recharge ensuite la liste par un nouvel appel \texttt{GET /api/candidats}.

Modifier un candidat. Le gestionnaire ouvre la fiche «~Modifier~» (données métier et, le cas échéant, affectation manuelle~: centre, établissement, salle, place). L'application envoie un \texttt{PUT /api/candidats/\{numeroInscription\}}. Le service valide le concours et l'affectation, met à jour PostgreSQL, puis la fiche affichée est actualisée.

Supprimer un candidat. Après confirmation, l'application envoie un \texttt{DELETE /api/candidats/\{numeroInscription\}}. Le service supprime l'enregistrement dans PostgreSQL et le candidat disparaît de la liste.

Réinitialiser la liste. Après confirmation, l'application envoie un \texttt{DELETE /api/candidats}. Le service supprime l'ensemble des candidats. La réponse \texttt{200 OK} indique le nombre d'éléments supprimés et l'écran affiche une liste vide.

\subsubsection{Diagramme de séquence --- gestion des concours}

Le diagramme de la figure~\ref{fig:seq-concours} décrit les opérations du gestionnaire sur la page «~Concours~». Cinq participants interviennent~: le gestionnaire, l'application web (React), la passerelle API, le service concours (Spring Boot) et PostgreSQL. Toutes les requêtes transitent par la passerelle, avec l'en-tête \texttt{Authorization: Bearer}. Un concours porte un numéro, une date et une heure d'examen, et au moins un centre.

\fullpagefig{seq-concours.png}{Diagramme de séquence --- gestion des concours}{fig:seq-concours}

Le scénario se déroule en quatre temps, correspondant aux cadres du diagramme.

Consulter la liste. À l'ouverture de la page, l'application web demande la liste par \texttt{GET /api/concours}. La passerelle transmet la requête au service concours, qui lit PostgreSQL (concours et centres affectés). La réponse \texttt{200 OK} ramène la liste, ensuite affichée à l'écran.

Ajouter un concours. Le gestionnaire saisit le nom, le numéro, la date et l'heure d'examen, ainsi qu'au moins un centre. L'application envoie un \texttt{POST /api/concours}. Le service vérifie l'unicité du numéro et l'existence des centres. Si le numéro est déjà utilisé, il répond \texttt{409 Conflict} et un message d'erreur s'affiche. Sinon, le concours et les affectations de centres sont enregistrés (\texttt{201 Created}) et l'application recharge la liste par un nouvel appel \texttt{GET /api/concours}.

Modifier un concours. Le gestionnaire ouvre la fiche «~Modifier~» et met à jour les champs ou les centres. L'application envoie un \texttt{PUT /api/concours/\{numeroConcours\}}. Le service valide les données, met à jour PostgreSQL, puis l'affichage est actualisé.

Supprimer un concours. Après confirmation, l'application envoie un \texttt{DELETE /api/concours/\{numeroConcours\}}. Le service supprime l'enregistrement dans PostgreSQL et le concours disparaît de la liste.

\subsubsection{Diagramme de séquence --- gestion des centres}

Le diagramme de la figure~\ref{fig:seq-centres} décrit les opérations du gestionnaire sur la page des lieux, pour le référentiel des centres (pôles géographiques, par exemple une ville d'examen). Cinq participants interviennent~: le gestionnaire, l'application web (React), la passerelle API, le service lieux (Spring Boot) et PostgreSQL. Toutes les requêtes transitent par la passerelle, avec l'en-tête \texttt{Authorization: Bearer}. Un centre est identifié par un nom unique. Les concours associés y apparaissent en lecture~; leur rattachement se gère depuis la page des concours.

\fullpagefig{seq-centres.png}{Diagramme de séquence --- gestion des centres}{fig:seq-centres}

Le scénario se déroule en quatre temps, correspondant aux cadres du diagramme.

Consulter les centres. À l'ouverture de la page, l'application web demande la liste par \texttt{GET /api/centres}. La passerelle transmet la requête au service lieux, qui exécute un \texttt{SELECT} sur PostgreSQL. La réponse \texttt{200 OK} ramène la liste, ensuite affichée à l'écran. La sélection d'un centre déclenche un \texttt{GET /api/centres/\{id\}}, qui ramène le détail, y compris les établissements rattachés.

Ajouter un centre. Le gestionnaire saisit uniquement le nom. L'application envoie un \texttt{POST /api/centres}. Le service vérifie l'unicité du nom. Si le nom est déjà utilisé, il répond \texttt{409 Conflict} et un message d'erreur s'affiche. Sinon, le centre est enregistré (\texttt{201 Created}), un identifiant \texttt{id\_centre} est généré, et la liste est mise à jour.

Modifier un centre. Le gestionnaire met à jour le nom. L'application envoie un \texttt{PUT /api/centres/\{id\}}. Le service contrôle à nouveau l'unicité, met à jour PostgreSQL, puis l'affichage est actualisé.

Supprimer un centre. Après confirmation, l'application envoie un \texttt{DELETE /api/centres/\{id\}}. Le service supprime le centre dans PostgreSQL~; la suppression entraîne, en cascade, les établissements et les salles rattachés. Le centre disparaît ensuite de la liste.

\subsubsection{Diagramme de séquence --- gestion des établissements}

Le diagramme de la figure~\ref{fig:seq-etablissements} décrit les opérations du gestionnaire sur les établissements d'un centre déjà ouvert. Cinq participants interviennent~: le gestionnaire, l'application web (React), la passerelle API, le service lieux (Spring Boot) et PostgreSQL. Toutes les requêtes transitent par la passerelle, avec l'en-tête \texttt{Authorization: Bearer}. Un établissement appartient toujours à un centre et porte un nom unique dans ce centre. Les concours y apparaissent en lecture, via les salles.

\fullpagefig{seq-etablissements.png}{Diagramme de séquence --- gestion des établissements}{fig:seq-etablissements}

Le scénario se déroule en quatre temps, correspondant aux cadres du diagramme.

Consulter les établissements. La liste ne fait pas l'objet d'un endpoint dédié~: elle est obtenue par \texttt{GET /api/centres/\{id\}}. La passerelle transmet la requête au service lieux, qui lit PostgreSQL et renvoie le centre avec ses établissements (noms et salles). La réponse \texttt{200 OK} permet d'afficher les établissements du centre ouvert.

Ajouter un établissement. Le gestionnaire saisit uniquement le nom. L'application envoie un \texttt{POST /api/centres/\{id\}/etablissements}. Le service vérifie l'unicité du nom dans ce centre. Si le nom est déjà utilisé, il répond \texttt{409 Conflict} et un message d'erreur s'affiche. Sinon, l'établissement est enregistré (\texttt{201 Created}) et un identifiant \texttt{id\_etablissement} est généré.

Modifier un établissement. Le gestionnaire met à jour le nom. L'application envoie un \texttt{PUT /api/etablissements/\{id\}}. Le service contrôle à nouveau l'unicité, met à jour PostgreSQL, puis l'affichage est actualisé.

Supprimer un établissement. Après confirmation, l'application envoie un \texttt{DELETE /api/etablissements/\{id\}}. Le service supprime l'établissement dans PostgreSQL~; la suppression entraîne, en cascade, les salles rattachées. L'établissement disparaît ensuite de la liste.

\subsubsection{Diagramme de séquence --- gestion des salles}

Le diagramme de la figure~\ref{fig:seq-salles} décrit les opérations du gestionnaire sur les salles d'un établissement, dans un centre déjà ouvert. Cinq participants interviennent~: le gestionnaire, l'application web (React), la passerelle API, le service lieux (Spring Boot) et PostgreSQL. Toutes les requêtes transitent par la passerelle, avec l'en-tête \texttt{Authorization: Bearer}. Une salle appartient toujours à un établissement, porte un nom unique dans cet établissement, une capacité, et au plus un concours. L'endpoint \texttt{GET /api/salles?numeroConcours=} existe par ailleurs pour la répartition et le tableau de bord, pas pour cette page.

\fullpagefig{seq-salles.png}{Diagramme de séquence --- gestion des salles}{fig:seq-salles}

Le scénario se déroule en quatre temps, correspondant aux cadres du diagramme.

Consulter les salles. La liste ne fait pas l'objet d'un endpoint dédié~: elle est obtenue par \texttt{GET /api/centres/\{id\}}. La passerelle transmet la requête au service lieux, qui lit PostgreSQL et renvoie le centre, ses établissements et les salles imbriquées (nom, nombre de places, concours). La réponse \texttt{200 OK} permet d'afficher les salles.

Ajouter une salle. Le gestionnaire saisit le nom, le nombre de places et éventuellement un concours. L'application envoie un \texttt{POST /api/etablissements/\{id\}/salles}. Le service vérifie l'unicité du nom dans cet établissement et l'existence du concours. Si le nom est déjà utilisé, il répond \texttt{409 Conflict} et un message d'erreur s'affiche. Sinon, la salle est enregistrée (\texttt{201 Created}) et un identifiant \texttt{id\_salle} est généré.

Modifier une salle. Le gestionnaire met à jour le nom, les places et/ou le concours. L'application envoie un \texttt{PUT /api/salles/\{idSalle\}}. Le service répète les contrôles d'unicité et d'existence, met à jour PostgreSQL, puis l'affichage est actualisé.

Supprimer une salle. Après confirmation, l'application envoie un \texttt{DELETE /api/salles/\{idSalle\}}. Le service supprime la salle dans PostgreSQL et celle-ci disparaît de la liste.

\subsubsection{Diagramme de séquence --- répartition automatique}

Le diagramme de la figure~\ref{fig:seq-repartition} décrit le déclenchement unique de la répartition automatique par le gestionnaire. Huit participants interviennent~: le gestionnaire, l'application web (React), la passerelle API, le service répartition, les services concours, lieux et candidats, et PostgreSQL. La requête du navigateur transite par la passerelle, avec l'en-tête \texttt{Authorization: Bearer}~; le service répartition relaie ce jeton vers les trois services métier. C'est une redistribution complète~: les candidats non placés voient leur affectation remise à zéro.

Le diagramme est large (huit lignes de vie)~: il est donc présenté en page paysage, afin de rester lisible.

\clearpage
\begin{landscape}
\begin{figure}[H]
  \centering
  \includegraphics[width=\linewidth,height=0.86\textheight,keepaspectratio]{seq-repartition.png}
  \caption{Diagramme de séquence --- répartition automatique}
  \label{fig:seq-repartition}
\end{figure}
\end{landscape}

Le scénario se déroule en quatre temps, correspondant aux cadres du diagramme.

Déclenchement. Un clic sur «~Lancer la répartition~» envoie un \texttt{POST /api/repartition/run} via la passerelle vers le service répartition.

Lecture des référentiels. Le service répartition interroge d'abord \texttt{GET /api/concours}. Pour chaque concours, il demande les salles par \texttt{GET /api/salles?numeroConcours=} (salles, centres et capacités), puis charge les candidats par \texttt{GET /api/candidats} (ville et concours).

Calcul du plan (en mémoire). Les candidats sont traités dans l'ordre du numéro d'inscription. Pour chacun, les centres du concours sont classés par proximité géographique (même ville, même région, puis distance). Si une salle a encore de la place, une affectation est réservée (centre, établissement, salle, numéro de place). Sinon, une alerte est enregistrée (capacité dépassée, aucun centre, concours inconnu ou ville non géolocalisée).

Écriture des affectations et de l'historique. Les affectations sont écrites en lot par \texttt{PATCH /api/candidats/affectations}. L'exécution (affectations et alertes) est historisée en PostgreSQL. Le service répond \texttt{201 Created} avec la synthèse~; l'application affiche alors le résultat au gestionnaire.

\subsubsection{Diagramme de séquence --- gestion des convocations}

Les diagrammes des figures~\ref{fig:seq-convocations}, \ref{fig:seq-convocations-pdf} et~\ref{fig:seq-convocations-envoi} décrivent les opérations du gestionnaire sur les convocations. Les participants comprennent le gestionnaire, l'application web (React), la passerelle API, le service convocations, Gmail (SMTP) et PostgreSQL. Sur la figure~\ref{fig:seq-convocations}, les appels vers les services candidats, concours et lieux sont regroupés sur la ligne de vie «~Services métier~», les chemins d'API restant indiqués sur les messages. Les requêtes du navigateur transitent par la passerelle, avec l'en-tête \texttt{Authorization: Bearer}~; le service convocations relaie ce jeton vers les trois services métier. Une convocation n'est pas enregistrée comme entité~: elle est assemblée à la lecture, uniquement si l'affectation est complète. PostgreSQL ne conserve que l'historique des envois.

Le scénario est présenté en trois figures, une par opération, afin que les messages restent lisibles. Les diagrammes les plus larges sont en page paysage.

\clearpage
\begin{landscape}
\begin{figure}[H]
  \centering
  \includegraphics[width=\linewidth,height=0.90\textheight,keepaspectratio]{seq-convocations-1.png}
  \caption{Diagramme de séquence --- gestion des convocations (1/3)~: consulter la liste}
  \label{fig:seq-convocations}
\end{figure}
\end{landscape}

Consulter les convocations. À l'ouverture de la page, l'application demande l'aperçu par \texttt{GET /api/convocations}. Le service convocations assemble les données auprès des services candidats, concours et lieux (\texttt{GET /api/candidats}, \texttt{GET /api/concours}, puis \texttt{GET /api/salles?numeroConcours=} pour chaque concours concerné) et ne retient que les candidats dont l'affectation est complète (centre, établissement, salle et place). Un second appel, \texttt{GET /api/convocations/envois}, lit en PostgreSQL l'historique des envois. La liste et l'historique s'affichent ensuite à l'écran.

\begin{figure}[H]
  \centering
  \includegraphics[width=\textwidth,height=0.52\textheight,keepaspectratio]{seq-convocations-2.png}
  \caption{Diagramme de séquence --- gestion des convocations (2/3)~: consulter un PDF}
  \label{fig:seq-convocations-pdf}
\end{figure}

Consulter un PDF. Un clic sur «~Voir le PDF~» déclenche \texttt{GET /api/convocations/\{numeroInscription\}/pdf}. Le service assemble la convocation, génère le fichier et le renvoie (\texttt{200 OK}). Le PDF n'est pas stocké.

\clearpage
\begin{landscape}
\begin{figure}[H]
  \centering
  \includegraphics[width=\linewidth,height=0.90\textheight,keepaspectratio]{seq-convocations-3.png}
  \caption{Diagramme de séquence --- gestion des convocations (3/3)~: envoyer toutes les convocations}
  \label{fig:seq-convocations-envoi}
\end{figure}
\end{landscape}

Envoyer toutes les convocations. L'envoi groupé se fait par \texttt{POST /api/convocations/envoyer}. Si Gmail n'est pas configuré, la réponse est \texttt{503}~; s'il n'y a aucun candidat affecté, \texttt{400}. Sinon, pour chaque convocation, un PDF est généré et envoyé par e-mail (pièce jointe)~; chaque tentative est historisée (\texttt{ENVOYE} ou \texttt{ECHEC}) et le bilan revient en \texttt{200 OK}. L'historique peut être vidé par \texttt{POST /api/convocations/envois/reinitialiser}.

\subsection{Diagramme de classes général}

Le diagramme de classes décrit la structure statique du système : classes métier, attributs (y compris les libellés dénormalisés), opérations et relations. Il complète les cas d'utilisation (qui fait quoi) et les diagrammes de séquence (comment les échanges se déroulent).

Le diagramme est présenté en une seule vue d'ensemble, en page paysage, afin de conserver les relations entre tous les domaines.

\clearpage
\begin{landscape}
\begin{figure}[H]
  \centering
  \includegraphics[width=\linewidth,height=0.90\textheight,keepaspectratio]{classe-general.png}
  \caption{Diagramme de classes général}
  \label{fig:classe-general}
\end{figure}
\end{landscape}

Le diagramme regroupe les classes par domaine. Un utilisateur porte un rôle \texttt{ADMINISTRATEUR} ou \texttt{GESTIONNAIRE} ; les opérations \texttt{creer}, \texttt{modifier}, \texttt{supprimer} et \texttt{consulter} concernent les comptes gestionnaires, tandis que \texttt{sAuthentifier} s'applique aux deux rôles. Un concours est identifié par son numéro ; il a un nom, une date et une heure d'examen, et il est composé d'au moins un centre (\texttt{ConcoursAffectationCentre}), avec copie du \texttt{nomCentre}. Les lieux forment une composition : un centre (nom unique) contient des établissements (nom unique dans le centre), eux-mêmes composés de salles (nom unique dans l'établissement, capacité, au plus un concours). Un candidat, créé uniquement par \texttt{importer} (fichier Excel), est identifié par son numéro d'inscription ; le CIN est unique. Il porte son identité complète, le concours (\texttt{numeroConcours} et \texttt{nomConcours} recopié), et, après répartition, éventuellement un siège (\texttt{idCentre}, \texttt{idEtablissement}, \texttt{idSalle}, \texttt{numeroPlace}). La répartition s'historise dans un \texttt{RepartitionRun} (\texttt{lancer}, \texttt{reinitialiserAffectations}), composé d'affectations et d'alertes (\texttt{CAPACITE\_DEPASSEE}, \texttt{AUCUN\_CENTRE\_DISPONIBLE}, \texttt{CONCOURS\_INCONNU}, \texttt{VILLE\_NON\_GEOLOCALISEE}). Chaque \texttt{RepartitionAffectation} recopie les libellés du candidat, du concours, du centre, de l'établissement et de la salle, afin que l'historique d'un run reste lisible si une source est renommée. Une convocation n'est pas persistée : elle est <<assemblée>> à la lecture si l'affectation est complète, avec les noms du candidat et des lieux (\texttt{consulter}, \texttt{telechargerPdf}, \texttt{envoyerToutes}). Seul l'envoi est enregistré (\texttt{ConvocationEnvoi}, statuts \texttt{ENVOYE} ou \texttt{ECHEC}), avec copie du nom et de l'e-mail. Les losanges pleins marquent la composition ; les flèches pointillées, une référence par identifiant ; la contrainte \{copie\} désigne un libellé dénormalisé.

\section{Conclusion}

Ce chapitre a présenté la conception de la plateforme de dématérialisation de la gestion des candidatures, selon une approche UML. Le diagramme de cas d'utilisation général a clarifié les fonctionnalités autour de deux acteurs : le gestionnaire, responsable de l'écriture métier, et l'administrateur, en lecture sur le métier et seul habilité à gérer les comptes gestionnaires. Les diagrammes de séquence ont décrit la dynamique des échanges entre l'application web, la passerelle API et les services concernés : authentification, candidatures (import Excel), concours, centres, établissements, salles, répartition automatique et convocations. Le diagramme de classes a formalisé la structure des entités métier et leurs relations (composition des lieux, références par identifiant, convocation assemblée à la lecture).

Cette conception fixe les règles essentielles du système. Les candidats sont créés uniquement par import. Un concours s'appuie sur au moins un centre. Les salles portent une capacité et au plus un concours. La répartition place les candidats selon la proximité géographique et la capacité des salles. Une convocation n'existe que si l'affectation est complète (centre, établissement, salle et place) ; elle n'est pas stockée, et seul l'historique des envois par e-mail est conservé.

Cette modélisation fournit une base cohérente pour la suite du rapport. Le chapitre suivant présente les outils utilisés et illustre l'application par des captures d'écran.

\end{onehalfspace}
